\documentclass[11pt]{article}

\usepackage{amsmath}
\usepackage{amssymb}
\usepackage{amsthm}
\usepackage[margin=1in]{geometry}

\newtheorem{definition}{Definition}[section]
\newtheorem{proposition}[definition]{Proposition}
\newtheorem{lemma}[definition]{Lemma}
\newtheorem{corollary}[definition]{Corollary}
\newtheorem{conjecture}[definition]{Conjecture}
\theoremstyle{remark}
\newtheorem{remark}[definition]{Remark}

\title{A counterexample to distributive-lattice layer-thickness rigidity}
\author{earthking11}
\date{13 September 2026}

\begin{document}
\maketitle

\begin{abstract}
We refute the conjecture that a finite distributive lattice is determined, up to
isomorphism, by its layer-count vector.  Since a finite distributive lattice is
isomorphic to the ideal lattice $J(P)$ of a finite poset $P$ (Birkhoff's theorem),
the layer-count vector is the rank-size vector of $J(P)$, namely the numbers of
order ideals of $P$ of each cardinality.  We exhibit non-isomorphic posets with
equal ideal-count vectors on $4$ elements, and in fact already on $4$ elements;
exhaustive enumeration shows that $4$ is the minimal size.  The natural reading of
the conjecture is therefore false.  We also record the corresponding collision on
$5$ elements.
\end{abstract}

\section{Setup and the conjecture}

Throughout, a \emph{poset} is a finite set $P$ with a partial order $\le$.  A
subset $I \subseteq P$ is an \emph{order ideal} (down-set) if $x \in I$ and
$y \le x$ imply $y \in I$.  The set $J(P)$ of order ideals, ordered by inclusion,
is a distributive lattice, and by Birkhoff's representation theorem every finite
distributive lattice arises this way, uniquely up to isomorphism: $J(P) \cong J(Q)$
if and only if $P \cong Q$.

\begin{definition}[Layer-count vector]
Let $L$ be a finite distributive lattice, and let $P$ be a poset with
$L \cong J(P)$.  The \emph{layer-count vector} of $L$ is
\[
  v(L) \;=\; \bigl(v_0, v_1, \dots, v_n\bigr), \qquad
  v_s \;=\; \#\{\, I \in J(P) : |I| = s \,\},
\]
where $n = |P|$.  Equivalently, $v_s$ is the number of elements of rank $s$ in
$L \cong J(P)$.
\end{definition}

The parameter name ``coatoum-layer thicknesses'' in the original statement is
garbled; the surrounding language (``layer-count vector'', ``Kruskal--Katona type
inequalities'', ``distributive lattice layer thickness rigidity'') describes the
rank-size statistics of the ideal lattice, which is the reading above.

\begin{conjecture}[Layer-thickness rigidity --- refuted]
\label{conj}
Finite distributive lattices with the same layer-count vector are isomorphic.
Equivalently, for finite posets $P, Q$: if $v(J(P)) = v(J(Q))$, then
$P \cong Q$.
\end{conjecture}

\begin{remark}
The conjecture is a statement about the map
$P \mapsto v(J(P))$.  It asserts that this invariant is complete on isomorphism
classes of finite posets.
\end{remark}

\section{A collision on $4$ elements}
\label{sec:4}

\begin{definition}[The two posets]
Let $P$ and $Q$ be the posets on $\{0,1,2,3\}$ whose strict order relations are
\[
  P:\quad 0 < 1 < 2, \quad 3 \text{ isolated};
  \qquad
  Q:\quad 0 < 2, \quad 0 < 3, \quad 1 < 2,
\]
each closed reflexively and transitively.  Explicitly, the nontrivial
comparabilities are
\[
  P:\; 0<1,\;0<2,\;1<2;
  \qquad
  Q:\; 0<2,\;0<3,\;1<2 .
\]
\end{definition}

\begin{proposition}[Equal layer-count vectors]
\label{prop:equal4}
$v(J(P)) = v(J(Q)) = (1,2,2,2,1)$.
\end{proposition}

\begin{proof}
We list the order ideals by cardinality.

\smallskip
\noindent\emph{Poset $P$ (chain $0<1<2$, isolated $3$).}
\[
\begin{array}{c|l}
s & \text{ideals of size } s \\ \hline
0 & \varnothing \\
1 & \{0\},\ \{3\} \\
2 & \{0,1\},\ \{0,3\} \\
3 & \{0,1,2\},\ \{0,1,3\} \\
4 & \{0,1,2,3\}
\end{array}
\]
giving $(1,2,2,2,1)$.

\smallskip
\noindent\emph{Poset $Q$ ($0<2$, $0<3$, $1<2$).}
\[
\begin{array}{c|l}
s & \text{ideals of size } s \\ \hline
0 & \varnothing \\
1 & \{0\},\ \{1\} \\
2 & \{0,1\},\ \{0,3\} \\
3 & \{0,1,2\},\ \{0,1,3\} \\
4 & \{0,1,2,3\}
\end{array}
\]
giving $(1,2,2,2,1)$ as well.
\end{proof}

\begin{lemma}[Non-isomorphism certificate]
\label{lem:noniso4}
$P \not\cong Q$.
\end{lemma}

\begin{proof}
The element $3$ of $P$ is isolated: it is comparable to no other element of $P$.
An order isomorphism preserves the property of being isolated (it is defined by the
order relation alone, and is transported by any bijection that preserves and
reflects $\le$).  On the other hand $Q$ has no isolated element: every element of
$Q$ is comparable to a distinct element ($0<2$, $0<3$, $1<2$, and dually
$2>0$, $3>0$, $2>1$).  Hence no bijection $\{0,1,2,3\} \to \{0,1,2,3\}$ is an
order isomorphism.

For completeness, one may also check directly: a direct computer enumeration of
all $4! = 24$ bijections confirms that none satisfies
$\bigl(a \le_P b \iff \sigma(a) \le_Q \sigma(b)\bigr)$ for all $a,b$.  This
computation is formalised in \texttt{lean4/Main.lean} as the theorem
\texttt{not\_isomorphic}, and reproduced in \texttt{reproduce.py}.
\end{proof}

\begin{corollary}
Conjecture~\ref{conj} is false.  The distributive lattices $J(P)$ and $J(Q)$ are
non-isomorphic yet have the same layer-count vector $(1,2,2,2,1)$.
\end{corollary}

\begin{proof}
By Birkhoff's theorem, $J(P) \cong J(Q)$ iff $P \cong Q$; Lemma~\ref{lem:noniso4}
gives $P \not\cong Q$, hence $J(P) \not\cong J(Q)$.  By
Proposition~\ref{prop:equal4} their layer-count vectors coincide.
\end{proof}

\section{A collision on $5$ elements}

The same phenomenon occurs with two posets whose layer-counts are spread less
evenly.

\begin{definition}
Let $A$ and $B$ be the posets on $\{0,1,2,3,4\}$ with strict relations
\[
  A:\; 0<1,\ 0<2,\ 0<3,\ 1<3,\ 2<3, \quad 4 \text{ isolated};
  \qquad
  B:\; 1<2,\ 2<3,\ 0<4 .
\]
\end{definition}

\begin{proposition}
\label{prop:equal5}
$v(J(A)) = v(J(B)) = (1,2,3,3,2,1)$.
\end{proposition}

\begin{proof}
Direct enumeration of the $2^5 = 32$ subsets of each poset and selection of the
down-closed ones.  For $A$ the ideal counts by size are $1,2,3,3,2,1$; for $B$ the
same.  The enumeration is carried out in \texttt{reproduce.py} and, up to the
sizes needed, in \texttt{lean4/Main.lean}.
\end{proof}

\begin{lemma}
\label{lem:noniso5}
$A \not\cong B$.
\end{lemma}

\begin{proof}
$A$ has an isolated element ($4$), while $B$ has none: $0<4$, $1<2<3$ make every
element of $B$ comparable to another.  Since the number of isolated elements is
an isomorphism invariant, $A \not\cong B$.  Equivalently, $A$ has $5$ nontrivial
comparabilities and $B$ has $3$ (after transitive closure), another isomorphism
invariant that distinguishes them.  A brute-force check of all $5! = 120$
bijections confirms this.
\end{proof}

\begin{corollary}
$J(A)$ and $J(B)$ are non-isomorphic finite distributive lattices with equal
layer-count vector $(1,2,3,3,2,1)$.
\end{corollary}

\section{Minimality: the smallest collision is on $4$ elements}

\begin{proposition}
\label{prop:minimal}
There is no pair of non-isomorphic posets on at most $3$ elements with equal
layer-count vectors.  The collision of Section~\ref{sec:4} on $4$ elements is
therefore minimal.
\end{proposition}

\begin{proof}
Exhaustive enumeration of all labelled reflexive, transitive, antisymmetric
relations on $n$ elements for $n = 1,2,3,4$, grouping by ideal-count vector and
testing every pair in a group for isomorphism by brute force over all $n!$
bijections, yields:
\[
\begin{array}{c|c|c|c}
n & \#\text{labelled posets} & \#\text{distinct vectors} & \#\text{non-isomorphic colliding pairs}\\ \hline
1 & 1   & 1  & 0 \\
2 & 3   & 2  & 0 \\
3 & 19  & 5  & 0 \\
4 & 219 & 15 & 576
\end{array}
\]
For $n \le 3$ every group of posets sharing an ideal-count vector is a single
isomorphism class, so the invariant is complete there.  For $n = 4$ there are
$576$ non-isomorphic colliding pairs, the pair $\{P,Q\}$ of Section~\ref{sec:4}
being one.  The enumeration is reproduced by \texttt{reproduce.py};
the collision on $4$ elements is formalised in \texttt{lean4/Main.lean}.
\end{proof}

\begin{remark}
The count of labelled posets $1, 3, 19, 219$ is the classical sequence of posets
on $n$ labelled points, a useful sanity check on the enumeration.
\end{remark}

\section{On the reading of ``layer-count vector''}

The refutation above targets the reading implicit in the statement's own language:
the layer-count vector records the number of elements of each rank of the
distributive lattice, equivalently the number of order ideals of each size in the
associated poset.  Under this reading the invariant is plainly not complete, as
the two examples show.

If instead ``layer-count vector'' were replaced by a richer invariant that already
determines the poset up to isomorphism --- for instance the isomorphism type of the
subposet of join-irreducible elements of the lattice, which by Birkhoff's theorem
recovers $P$ --- then the assertion ``same invariant $\Rightarrow$ isomorphic''
would be trivially true.  That is a different, tautological statement.  The
conjecture as written, with its reference to layer counts and Kruskal--Katona type
inequalities on those counts, is the rank-statistics reading, and it is false.

\end{document}
